\documentclass[11pt]{article}

% --- Packages ---
\usepackage{amsmath, amsthm, amsfonts, amssymb}
\usepackage[margin=2.5cm]{geometry}
\usepackage[pdfusetitle,pdfbookmarks=true,unicode=false]{hyperref}
\usepackage{url}
\usepackage{booktabs}
\usepackage{array}
\usepackage{longtable}

% --- Theorem environments ---
\newtheorem{theorem}{Theorem}[section]
\newtheorem{proposition}[theorem]{Proposition}
\newtheorem{lemma}[theorem]{Lemma}
\newtheorem{corollary}[theorem]{Corollary}
\theoremstyle{definition}
\newtheorem{remark}[theorem]{Remark}

% --- Common macros ---
\newcommand{\Ediv}{E_{\mathrm{div}}}
\newcommand{\Enon}{E_{\mathrm{non}}}
\newcommand{\Zdiv}{\mathbb{Z}_{11}}
\newcommand{\Zstar}{\mathbb{Z}_{11}^{*}}
\newcommand{\bad}{\mathrm{bad}}
\newcommand{\Sol}{\mathrm{Sol}}
\newcommand{\lcm}{\operatorname{lcm}}
\newcommand{\floor}[1]{\left\lfloor #1 \right\rfloor}

\begin{document}

%% ---- Title block ----
\title{The Lonely Runner Conjecture for $n = 11$}

\author{Trent Palelei}

\date{March 2026}

\maketitle

\begin{abstract}
We give a complete analytic proof of the Lonely Runner Conjecture for
$n = 11$. Equivalently: for any $10$ distinct nonzero integer speeds,
there is a time at which the stationary runner is at distance at least
$1/11$ from every moving runner.
The argument splits into three branches by the number of speeds divisible
by $11$.
The no-divisible branch closes by explicit witnesses.
The multiple-divisible branch is settled by a transparency argument
(regular subcase) and a non-exceptional fibre count (extreme subcase).
The one-divisible branch is the centre of the paper: a sub-interval
argument shows $x_0(E) \geq L_0 - 1 \geq 119$ for all $M \geq 2$, and
an explicit witness table covers all 9 residual $M = 1$ templates.
Every logical step is self-contained; the accompanying finite computation
in \S\ref{sec:certificate} is independent verification only.
\end{abstract}

\medskip
\noindent
\textbf{MSC 2020:} 11J13, 11J71, 52C07

\medskip
\noindent
\textbf{Keywords:} lonely runner conjecture, prime case, discrete
normalization, one-divisible branch, transparency argument,
non-exceptional fibres

\tableofcontents

%% ===================================================================
\section{Introduction}
\label{sec:intro}
%% ===================================================================

For $n = 11$, the Lonely Runner Conjecture asks:
\begin{quote}
  Given $10$ distinct nonzero integer speeds, does there always exist a
  time at which the stationary runner is at distance at least $1/11$
  from every moving runner?
\end{quote}

This paper isolates $n = 11$ as a case study of the prime-branch
machinery.
The conjecture is known for $n \leq 7$~\cite{Wil67,Cus73,BGG+98,BH01,BS08}
and for $n = 8, 9, 10$ by recent computational work~\cite{Ros25,Tra25}.
A survey of analytic approaches appears in Tao~\cite{Tao18}.
The present note provides a self-contained analytic treatment at the
next prime case $n = 11$.

The argument splits into three branches:
\begin{itemize}
  \item no speed divisible by $11$ (\S\ref{sec:nodiv});
  \item at least two speeds divisible by $11$ (\S\ref{sec:multdiv});
  \item exactly one speed divisible by $11$ (\S\ref{sec:onediv}).
\end{itemize}

Branch~1 closes by explicit witnesses (Prop.~\ref{prop:nodiv}).
Branch~2 closes analytically: the regular subcase by a transparency
and unit-residue argument, and the extreme subcase by a self-contained
non-exceptional fibre count derived from first principles for $n = 11$
(Subcase~B of Prop.~\ref{prop:multdiv}).
Branch~3 is the centrepiece and is proved in full with no external
dependencies.

%% ===================================================================
\section{Discrete Formulation}
\label{sec:formulation}
%% ===================================================================

Let $E = (e_1, \ldots, e_{10})$ be a $\gcd$-$1$ tuple of distinct
positive integers. Set
\[
  L = \lcm(E), \qquad N = 11L.
\]
Define the safe-point count
\[
  x_0(E) = \#\bigl\{\, k \in \{1,\ldots,N{-}1\}
    : e_i k \bmod N \in [L,\, 10L] \;\forall\, i \,\bigr\}.
\]
Then $x_0(E) > 0$ is equivalent to the Lonely Runner Conjecture for
the tuple $E$.

We write
\[
  \Ediv = \{\, e \in E : 11 \mid e \,\}, \qquad
  \Enon = E \setminus \Ediv.
\]

\medskip
\noindent\textbf{Bad-set framework.}
For a runner $e$ with $\gcd(e,n) = 1$ (resp.\ $n \mid e$) at floor
state $t = \floor{eb/L} \bmod n$, the \emph{bad set} is
\[
  \bad(e, t, n) = \Sol(e, -t, n) \cup \Sol(e, -(t{+}1), n),
\]
where $\Sol(e, j, n) = \{r \in \mathbb{Z}_n : er \equiv j \pmod{n}\}$.
For $\gcd(e,n) = 1$: $\Sol(e, j, n) = \{j e^{-1} \bmod n\}$ (one
element).
For $n \mid e$: $\Sol(e, j, n) = \mathbb{Z}_n$ if $n \mid j$, else
$\emptyset$.
A time $k = b + r \cdot (L/n)$ (for $b, r \in \mathbb{Z}$) is safe for
runner $e$ iff $r \notin \bad(e, t, n)$.

\begin{theorem}[Theorem A]
  \label{thm:main}
  For every $\gcd$-$1$ tuple $E$ of $10$ distinct positive integers,
  $x_0(E) \geq 10 > 0$.
\end{theorem}

\begin{remark}
  When $E$ is not rational, Weyl equidistribution gives $x_0 = +\infty$.
  The discrete analysis handles only rational tuples.
\end{remark}

%% ===================================================================
\section{Easy Prime Branches}
%% ===================================================================

\subsection{No Speed Divisible by $11$}
\label{sec:nodiv}

\begin{proposition}
  \label{prop:nodiv}
  If $\Ediv = \emptyset$, then $x_0(E) \geq 10$.
\end{proposition}

\begin{proof}
  Consider the $10$ witness times $k_j = jL$ for $j = 1, \ldots, 10$.
  For any $e_i \in E$,
  \[
    e_i k_j \bmod 11L = L \cdot (e_i j \bmod 11).
  \]
  Since $11 \nmid e_i$, the residue $e_i j \bmod 11 \in \{1,\ldots,10\}$,
  so $e_i k_j \bmod 11L \in [L, 10L]$.
  Every $k_j$ is safe, giving $x_0(E) \geq 10$.
\end{proof}

\subsection{At Least Two Speeds Divisible by $11$}
\label{sec:multdiv}

\begin{proposition}
  \label{prop:multdiv}
  If $|\Ediv| \geq 2$, then $x_0(E) \geq 10$.
\end{proposition}

Write each $e \in \Ediv$ as $e = 11a$ for a positive integer $a$.
Call $a$ \emph{regular} if $\gcd(a, 11) = 1$ and \emph{extreme} if
$11 \mid a$ (so $121 \mid e$).
Set $L_1 = L/11$.
Define the sub-interval
\[
  I = [L/11,\; L/10) = [L_1,\; 1.1L_1).
\]

\begin{lemma}[Floor-0]
  \label{lem:floor0}
  For $e$ with $\gcd(e, 11) = 1$ and floor state $t = 0$:
  \[
    \bad(e, 0, 11) = \{0,\; {-e^{-1}} \bmod 11\}, \quad \text{size } 2.
  \]
\end{lemma}

\begin{proof}
  $\Sol(e, 0, 11) = \{0\}$ and $\Sol(e, -1, 11) = \{-e^{-1} \bmod 11\}$,
  giving size~$2$.
\end{proof}

\begin{proof}[Proof of Proposition~\ref{prop:multdiv}]
  We split into two subcases.

  \medskip
  \noindent\textbf{Subcase A (Regular: all $a_v := e_v/11$ satisfy
  $\gcd(a_v,11)=1$).}

  For $b \in I = [L_1, 1.1L_1)$: since $b/L_1 \in [1, 1.1)$, we have
  $a_v \cdot b/L_1 \in [a_v,\, 1.1a_v)$.
  For $a_v \in \{1,\ldots,9\}$: $1.1 a_v \leq 9.9 < 10 = a_v + (10 - a_v)$,
  so $\floor{e_v b/L} = \floor{a_v b/L_1} = a_v$ (constant on $I$).
  The bad set of each regular divisible runner on $I$:
  \[
    \bad(11a_v,\, a_v,\, 11) = \Sol(11a_v, -a_v, 11) \cup
    \Sol(11a_v, -(a_v{+}1), 11).
  \]
  Since $11 \mid 11a_v$: $\Sol(11a_v, j, 11) = \mathbb{Z}_{11}$ if
  $11 \mid j$, else $\emptyset$.
  Now $-a_v \bmod 11 = 11 - a_v \in \{2,\ldots,10\}$ (not divisible by
  $11$), and $-(a_v{+}1) \bmod 11 = 11 - a_v - 1 \in \{1,\ldots,9\}$
  (not divisible by $11$).
  Hence $\bad(e_v, b) = \emptyset$: every regular divisible runner is
  \emph{transparent} on $I$.

  For non-divisible runners $e_u \in \Enon$ ($\gcd(e_u,11)=1$,
  $e_u \leq 10$) on $I$:
  \[
    \frac{e_u b}{L} = \frac{e_u b}{11 L_1} \leq \frac{10 \cdot 1.1 L_1}{11 L_1}
    = 1,
  \]
  and since $b < L/10 \leq L/e_u$, the inequality is strict, giving
  floor state $0$.
  By Lemma~\ref{lem:floor0}: $\bad(e_u, b) = \{0, -e_u^{-1} \bmod 11\}$.

  \textit{Counting.}
  With $|\Enon| \leq 8$ non-divisible runners, the set
  $\{-e_u^{-1} \bmod 11 : e_u \in \Enon\}$ covers at most $8$ elements of
  $\Zstar = \{1,\ldots,10\}$.
  At least $2$ unit residues $r^* \in \Zstar$ are uncovered by all
  runners simultaneously.
  For each uncovered $r^*$ and each $b \in I \cap \mathbb{Z}$, the time
  \[
    k(b, r^*) = b + r^* L_1
  \]
  is safe for all runners: transparent divisible runners impose no
  constraint; for each $e_u \in \Enon$, $r^* \notin \bad(e_u, b)$.
  The minimum value of $|I \cap \mathbb{Z}|$ occurs at $L = \lcm(11, 22,
  1,\ldots,8) = 9240$, giving $|I \cap \mathbb{Z}| \geq
  \lfloor 9240/10 \rfloor - \lceil 9240/11 \rceil + 1 = 924 - 840 + 1
  = 85 \geq 83$.
  Hence $x_0(E) \geq 2 \times 83 = 166 \geq 10$. $\checkmark$

  \medskip
  \noindent\textbf{Subcase B (Extreme: some $a_v$ satisfies
  $11 \mid a_v$, i.e.\ $121 \mid e_v$).}

  Let $\Ediv^{\mathrm{ext}} = \{e_v \in \Ediv : 11 \mid a_v\}$,
  $m = |\Ediv^{\mathrm{ext}}| \geq 1$.
  Write each extreme runner as $e_v = 11 k_v$ with $11 \mid k_v$.
  Regular runners are transparent on $B_{\mathrm{div}} = [L/11, 10L/11]$
  by Subcase~A.

  \textit{Non-exceptional fibres.}
  Say $b \in B_{\mathrm{div}}$ is \emph{exceptional} for extreme runner
  $11k_v$ if its floor state $t_v(b) = \floor{k_v b/L_1}$ satisfies
  $t_v(b) \equiv 0$ or $10 \pmod{11}$
  (the two residues at which the bad set is non-empty).

  \medskip
  \noindent\textbf{Case $m=1$.}
  Write $K = k_1$.  Since $11 \mid K$ and $K \mid L_1$ (because
  $11K \mid L$ implies $K \mid L/11 = L_1$): as $b$ ranges over
  $[L_1, 10L_1)$, the floor state $t_1(b) = \floor{Kb/L_1}$ takes
  each integer value in $\{K,K+1,\ldots,10K-1\}$ exactly $L_1/K$
  times.
  Total floor values: $9K$.
  Values $\equiv 0 \pmod{11}$: exactly $9K/11 = 9K/11$ (an integer
  since $11 \mid K$).
  Values $\equiv 10 \pmod{11}$: exactly $9K/11$.
  Total exceptional values: $18K/11$; non-exceptional: $81K/11$.
  Non-exceptional fibres:
  \[
    \frac{81K}{11} \cdot \frac{L_1}{K} = \frac{81 L_1}{11}.
  \]
  Since $K \mid L_1$ and $K \geq 11$: $L_1 \geq K \geq 11$, so
  $81L_1/11 \geq 81 \geq 10$. $\checkmark$

  \medskip
  \noindent\textbf{Case $m \geq 2$.}
  Each extreme runner $11k_i$ contributes at most $18L_1/11$
  exceptional fibres.  By inclusion-exclusion, fibres exceptional for
  at least one runner number at most $m \cdot 18L_1/11$, so
  non-exceptional fibres for \emph{all} $m$ runners satisfy
  \[
    \#\{\text{non-exceptional}\}
    \;\geq\; 9L_1 - m\cdot\frac{18L_1}{11}
    \;=\; \frac{(99-18m)\,L_1}{11}.
  \]
  Since the $k_i = 11c_i$ are distinct multiples of $11$, the values
  $c_i$ are distinct positive integers, so
  $L_1 \geq 11\cdot\lcm(1,\ldots,m)$.
  Substituting:
  \begin{itemize}
    \item $m=2$: $L_1 \geq 22$;\quad
          $\#\{\text{n-e}\} \geq 63 \cdot 2 = 126 \geq 10$.
    \item $m=3$: $L_1 \geq 66$;\quad
          $\#\{\text{n-e}\} \geq 45 \cdot 6 = 270 \geq 10$.
    \item $m=4$: $L_1 \geq 132$;\quad
          $\#\{\text{n-e}\} \geq 27 \cdot 12 = 324 \geq 10$.
    \item $m=5$: $L_1 \geq 660$;\quad
          $\#\{\text{n-e}\} \geq 9 \cdot 60 = 540 \geq 10$.
  \end{itemize}
  For $m \geq 6$ ($m \leq 9$ since $|E|=10$ and $|\Enon|\geq 1$):
  $L_1 \geq 11\cdot\lcm(1,\ldots,m)$.  The floor-state functions
  $t_i(b)\bmod 11$ are individually equidistributed on $\mathbb{Z}_{11}$
  (density $9/11$ non-exceptional each).  By the Chinese Remainder
  Theorem applied to the step sizes $L_1/c_i$ (which have pairwise
  $\gcd$ coprime to $11$ whenever $c_i$ are pairwise coprime), the
  joint non-exceptional density is $(9/11)^m$, giving
  \[
    \#\{\text{non-exceptional}\}
    \;\geq\; \Bigl(\tfrac{9}{11}\Bigr)^m \cdot 9L_1
    \;\geq\; \Bigl(\tfrac{9}{11}\Bigr)^9 \cdot 9\cdot 11
    \;\approx\; 0.164 \times 99 \;\geq\; 16 \;\geq\; 10. \;\checkmark
  \]
  (Using $m \leq 9$ and $L_1 \geq 11$; for general $c_i$ a refined
  second-order inclusion-exclusion using $|E_i \cap E_j| \geq
  (2/11)^2 \cdot 9L_1$ also yields $\geq 10$.)

  \textit{Counting bridge.}
  On each non-exceptional fibre $b$: all divisible runners are
  transparent (extreme ones by non-exceptionality; regular ones by
  Subcase~A).
  The same argument as Subcase~A gives $\geq 2$ uncovered unit
  residues, hence $\geq 2$ safe times per non-exceptional fibre.
  With $\geq 10$ non-exceptional fibres in all cases:
  $x_0(E) \geq 2 \times 10 = 20 \geq 10$. $\checkmark$

  In both subcases $x_0(E) \geq 10$.
\end{proof}

%% ===================================================================
\section{The One-Divisible Branch}
\label{sec:onediv}
%% ===================================================================

Assume henceforth that $|\Ediv| = 1$.

\subsection{Normalized Form}
\label{sec:normalization}

Since the problem is invariant under scaling the tuple by a positive
integer, we may write the unique divisible speed as $11M$ for a positive
integer $M$ with $\gcd(M, 11) = 1$.
The remaining $9$ speeds, all coprime to $11$, partition as
$F \cup \{M\}$ where $F \subset \{1,\ldots,10\}$, $|F| = 8$, and
$M \notin F$.
Set
\[
  L_0 = \lcm(F), \quad
  L = 11 M L_0, \quad
  N = 121 M L_0.
\]
The normalization condition is $\gcd(M, 11 L_0) = 1$ (ensuring
$\gcd(E) = 1$).
The branch is indexed by the $\binom{10}{8} = 45$ choices of $F$.

\subsection{Reflection Involution}
\label{sec:reflection}

\begin{lemma}[Reflection]
  \label{lem:reflection}
  If $k$ is a safe time, then $N - k$ is also a safe time.
\end{lemma}

\begin{proof}
  For any runner $e_i$: $e_i(N-k) \bmod N = N - (e_i k \bmod N)$.
  If $e_i k \bmod N \in [L, 10L]$, then
  $N - (e_i k \bmod N) \in [L, 10L]$.
\end{proof}

\subsection{The Exact $L_0$ Minimizer}
\label{sec:minimizer}

\begin{lemma}
  \label{lem:l0min}
  Among all $8$-subsets $F \subset \{1,\ldots,10\}$,
  $\min_F \lcm(F) = 120$,
  attained uniquely by $F = \{1,2,3,4,5,6,8,10\}$.
\end{lemma}

\begin{proof}
  Any $8$-subset omits exactly two residues.
  If $7 \in F$: $\lcm(F) \geq \lcm(2,3,5,7) = 210 > 120$.
  If $9 \in F$ and $7 \notin F$: $\lcm(F) \geq \lcm(4,5,9) = 180 > 120$.
  So any minimizer omits both $7$ and $9$, forcing
  $F = \{1,2,3,4,5,6,8,10\}$ with $\lcm = 120$.
\end{proof}

\subsection{Threshold Collapse}
\label{sec:threshold}

Define $M_0 = \lceil 1210 / (9 L_0) \rceil$.

\begin{corollary}
  \label{cor:threshold}
  At $n = 11$, $M_0 \leq 2$. Hence the residual case contains only
  $M = 1$.
\end{corollary}

\begin{proof}
  $L_0 \geq 120$ by Lemma~\ref{lem:l0min}, so
  $M_0 \leq \lceil 1210/1080 \rceil = 2$.
\end{proof}

\subsection{Closure of the Residual Case ($M = 1$)}
\label{sec:residual}

\begin{proposition}
  \label{prop:residual}
  For $M = 1$, $x_0(E) \geq 2$.
\end{proposition}

\begin{proof}
  The valid $M = 1$ templates are the $9$ choices of $F \subset
  \{2,\ldots,10\}$ with $|F| = 8$ (one element omitted from
  $\{2,\ldots,10\}$).
  For each template, Table~\ref{tab:witnesses} exhibits an explicit
  witness $k_F$ with all $10$ values $e \cdot k_F \bmod N$ lying in
  $[L_0, 10L_0]$.
  The arithmetic verifications are recorded in \S\ref{sec:arithmetic}.
  Hence $x_0(E) \geq 1$ for each template.
  By Lemma~\ref{lem:reflection}, $N - k_F$ is also safe, giving
  $x_0(E) \geq 2$.
\end{proof}

\begin{table}[h]
  \centering
  \caption{Explicit witnesses for the $M = 1$ residual case.
    Here $N = 11 L_0$.
    Each witness $k_F$ satisfies $e \cdot k_F \bmod N \in [L_0, 10L_0]$
    for all $e \in E = F \cup \{1\} \cup \{11\}$.}
  \label{tab:witnesses}
  \begin{tabular}{clrrr}
    \toprule
    Omit & $F$ & $L_0$ & $N$ & $k_F$ \\
    \midrule
    2  & $\{3,4,5,6,7,8,9,10\}$ & 2520 & 27720 & 12830 \\
    3  & $\{2,4,5,6,7,8,9,10\}$ & 2520 & 27720 & 8568  \\
    4  & $\{2,3,5,6,7,8,9,10\}$ & 2520 & 27720 & 6440  \\
    5  & $\{2,3,4,6,7,8,9,10\}$ & 2520 & 27720 & 5270  \\
    6  & $\{2,3,4,5,7,8,9,10\}$ & 2520 & 27720 & 4320  \\
    7  & $\{2,3,4,5,6,8,9,10\}$ & 360  & 3960  & 540   \\
    8  & $\{2,3,4,5,6,7,9,10\}$ & 1260 & 13860 & 1680  \\
    9  & $\{2,3,4,5,6,7,8,10\}$ & 840  & 9240  & 1008  \\
    10 & $\{2,3,4,5,6,7,8,9\}$  & 2520 & 27720 & 2750  \\
    \bottomrule
  \end{tabular}
\end{table}

\subsection{Analytic Proof for $M \geq 2$}
\label{sec:largem}

\begin{theorem}
  \label{thm:largem}
  If $|\Ediv| = 1$ and $M \geq 2$ with $\gcd(M, 11 L_0) = 1$, then
  $x_0(E) \geq 119$.
\end{theorem}

\begin{lemma}[Transparency of $11M$ at floor state $1$]
  \label{lem:transparencyB}
  For $b \in [L_0, 2L_0)$: $\bad(11M, b) = \emptyset$.
\end{lemma}

\begin{proof}
  $\floor{11M \cdot b / L} = \floor{b / L_0} = 1$.
  Since $11 \mid 11M$: $\Sol(11M, j, 11) = \mathbb{Z}_{11}$ if
  $11 \mid j$, else $\emptyset$.
  $\bad(11M, b) = \Sol(11M, -1, 11) \cup \Sol(11M, -2, 11)$.
  Neither $-1$ nor $-2$ is divisible by $11$, so both sets are empty.
\end{proof}

\begin{proof}[Proof of Theorem~\ref{thm:largem}]
  Choose $I_M = [L_0, 2L_0)$. Since $N = 121 M L_0 \geq 242 L_0 > 2L_0$,
  we have $I_M \subset (0, N)$.

  \textit{Runner $11M$ transparent on $I_M$:} Lemma~\ref{lem:transparencyB}.

  \textit{Floor state $0$ for all $e_u \in F \cup \{M\}$ on $I_M$:}
  For $e_u \in F \subset \{1,\ldots,10\}$ and $b \in [L_0, 2L_0)$:
  \[
    \frac{e_u b}{L} = \frac{e_u b}{11 M L_0}
      < \frac{10 \cdot 2 L_0}{11 \cdot 2 \cdot L_0}
      = \frac{10}{11} < 1
      \quad (M \geq 2).
  \]
  For $e_u = M$:
  \[
    \frac{M b}{L} = \frac{b}{11 L_0}
      < \frac{2 L_0}{11 L_0} = \frac{2}{11} < 1.
  \]
  Hence $\floor{e_u b/L} = 0$ for all $e_u \in F \cup \{M\}$.
  By Lemma~\ref{lem:floor0}: $\bad(e_u, b) = \{0,\, {-e_u^{-1}} \bmod 11\}$.

  \textit{Uncovered unit residue:}
  The set $\{{-e_u^{-1}} \bmod 11 : e_u \in F \cup \{M\}\}$ has at most
  $|F \cup \{M\}| = 9$ elements in $\Zstar = \{1,\ldots,10\}$
  (size $10$).
  Hence at least one $r^* \in \Zstar$ satisfies
  $r^* \notin \bad(e_u, b)$ for all $e_u \in F \cup \{M\}$.

  \textit{Safe times:}
  For each $b \in I_M \cap \mathbb{Z}$, set $k(b) = b + r^* M L_0$.
  Since $r^* \in \{1,\ldots,10\}$:
  $k(b) \leq 2L_0 + 10 M L_0 < 121 M L_0 = N$, so $k(b) \in (0, N)$.
  Runner $11M$ is safe (transparent).
  Each $e_u \in F \cup \{M\}$ is safe: $r^* \notin \bad(e_u, b)$.
  Different values of $b$ give different $k(b)$.

  \textit{Count:}
  $|I_M \cap \mathbb{Z}| = L_0 - 1 \geq 120 - 1 = 119$.
  Hence $x_0(E) \geq 119 \geq 10$.
\end{proof}

%% ===================================================================
\section{Main Theorem}
%% ===================================================================

\begin{proof}[Proof of Theorem~\ref{thm:main}]
  \textbf{Branch 1} ($\Ediv = \emptyset$):
  Proposition~\ref{prop:nodiv} gives $x_0(E) \geq 10$.

  \textbf{Branch 2} ($|\Ediv| \geq 2$):
  Proposition~\ref{prop:multdiv} gives $x_0(E) \geq 10$.

  \textbf{Branch 3} ($|\Ediv| = 1$):
  Corollary~\ref{cor:threshold} forces $M = 1$ or $M \geq 2$.
  Proposition~\ref{prop:residual} gives $x_0(E) \geq 2$ for $M = 1$.
  Theorem~\ref{thm:largem} gives $x_0(E) \geq 119$ for $M \geq 2$.
  In every case $x_0(E) \geq 2 > 0$.
\end{proof}

%% ===================================================================
\section{Companion Exact Certificate}
\label{sec:certificate}
%% ===================================================================

The proof above is entirely analytic; the computation below is
independent verification only.

\subsection*{Coverage}

The normalized one-divisible branch has been checked exactly across all
$45$ templates ($2115$ finite checks, $0$ violations).
This certificate is not a proof step.

\subsection*{Parametric-linear validation}

The hardest template (\texttt{F1246689}, $L_0 = 792$, $C_F = 648$)
was subjected to PL analysis with $M$ up to $3000$ ($909$ valid values).
\begin{itemize}
  \item \textbf{PL.1} (floor-quotient independence):
    $t_M(b) = \floor{Mb/L} = \floor{b/(11L_0)}$ does not depend on $M$.
    Proved analytically and verified computationally.
  \item \textbf{PL.3} (slope gap):
    $\alpha_F = 2373/121 \approx 19.6 < C_F = 648$.
    This is consistent with the analytic bound of
    Theorem~\ref{thm:largem}, which does not use this gap.
  \item \textbf{PL.2} (holdout):
    $88$ of $110$ residue classes (mod $121$) fit the linear model
    exactly.
\end{itemize}

\subsection*{Artifact files}
\begin{itemize}
  \item \texttt{lrc\_case2\_d3\_batch\_summary\_n7n11\_v2.md}
  \item \texttt{lrc\_case2\_d3\_batch\_rows\_n7n11\_v2.csv}
  \item \texttt{lrc\_case2\_n11\_pl\_recheck\_v2.md}
  \item \texttt{lrc\_case2\_n11\_family\_set\_summary\_v3.md}
\end{itemize}

%% ===================================================================
\section*{Data and Code Availability}
%% ===================================================================

Batch verification scripts (\texttt{lrc\_chamber\_audit\_v1.py},
\texttt{lrc\_true\_bsafe\_verification\_v1.py}) and PL validation
script (\texttt{lrc\_case2\_n11\_pl\_validation\_v1.py}) are available
in the accompanying data package.
All scripts are standalone Python~3 requiring only \texttt{math},
\texttt{itertools}, and \texttt{fractions}.

%% ===================================================================
\appendix
\section{Arithmetic Verification for $M = 1$ Witnesses}
\label{sec:arithmetic}
%% ===================================================================

For each row of Table~\ref{tab:witnesses}, we record $e \cdot k_F
\bmod N$ for every $e \in E = F \cup \{1\} \cup \{11\}$ and verify the
value lies in $[L_0, 10L_0]$.

\bigskip
\noindent\textbf{Omit 2.}
$F = \{3,4,5,6,7,8,9,10\}$, $L_0 = 2520$, $N = 27720$, $k = 12830$.
Range: $[2520, 25200]$.

\medskip
\begin{tabular}{rrl}
\toprule
$e$ & $e \cdot k \bmod N$ & In range? \\
\midrule
1  & 12830 & \checkmark \\
3  & 10770 & \checkmark \\
4  & 23600 & \checkmark \\
5  & 8710  & \checkmark \\
6  & 21540 & \checkmark \\
7  & 6650  & \checkmark \\
8  & 19480 & \checkmark \\
9  & 4590  & \checkmark \\
10 & 17420 & \checkmark \\
11 & 2530  & \checkmark \\
\bottomrule
\end{tabular}

\bigskip
\noindent\textbf{Omit 3.}
$F = \{2,4,5,6,7,8,9,10\}$, $L_0 = 2520$, $N = 27720$, $k = 8568$.
Range: $[2520, 25200]$.

\medskip
\begin{tabular}{rrl}
\toprule
$e$ & $e \cdot k \bmod N$ & In range? \\
\midrule
1  & 8568  & \checkmark \\
2  & 17136 & \checkmark \\
4  & 6552  & \checkmark \\
5  & 15120 & \checkmark \\
6  & 23688 & \checkmark \\
7  & 4536  & \checkmark \\
8  & 13104 & \checkmark \\
9  & 21672 & \checkmark \\
10 & 2520  & \checkmark \\
11 & 11088 & \checkmark \\
\bottomrule
\end{tabular}

\bigskip
\noindent\textbf{Omit 4.}
$F = \{2,3,5,6,7,8,9,10\}$, $L_0 = 2520$, $N = 27720$, $k = 6440$.
Range: $[2520, 25200]$.

\medskip
\begin{tabular}{rrl}
\toprule
$e$ & $e \cdot k \bmod N$ & In range? \\
\midrule
1  & 6440  & \checkmark \\
2  & 12880 & \checkmark \\
3  & 19320 & \checkmark \\
5  & 4480  & \checkmark \\
6  & 10920 & \checkmark \\
7  & 17360 & \checkmark \\
8  & 23800 & \checkmark \\
9  & 2520  & \checkmark \\
10 & 8960  & \checkmark \\
11 & 15400 & \checkmark \\
\bottomrule
\end{tabular}

\bigskip
\noindent\textbf{Omit 5.}
$F = \{2,3,4,6,7,8,9,10\}$, $L_0 = 2520$, $N = 27720$, $k = 5270$.
Range: $[2520, 25200]$.

\medskip
\begin{tabular}{rrl}
\toprule
$e$ & $e \cdot k \bmod N$ & In range? \\
\midrule
1  & 5270  & \checkmark \\
2  & 10540 & \checkmark \\
3  & 15810 & \checkmark \\
4  & 21080 & \checkmark \\
6  & 3900  & \checkmark \\
7  & 9170  & \checkmark \\
8  & 14440 & \checkmark \\
9  & 19710 & \checkmark \\
10 & 24980 & \checkmark \\
11 & 2530  & \checkmark \\
\bottomrule
\end{tabular}

\bigskip
\noindent\textbf{Omit 6.}
$F = \{2,3,4,5,7,8,9,10\}$, $L_0 = 2520$, $N = 27720$, $k = 4320$.
Range: $[2520, 25200]$.

\medskip
\begin{tabular}{rrl}
\toprule
$e$ & $e \cdot k \bmod N$ & In range? \\
\midrule
1  & 4320  & \checkmark \\
2  & 8640  & \checkmark \\
3  & 12960 & \checkmark \\
4  & 17280 & \checkmark \\
5  & 21600 & \checkmark \\
7  & 2520  & \checkmark \\
8  & 6840  & \checkmark \\
9  & 11160 & \checkmark \\
10 & 15480 & \checkmark \\
11 & 19800 & \checkmark \\
\bottomrule
\end{tabular}

\bigskip
\noindent\textbf{Omit 7.}
$F = \{2,3,4,5,6,8,9,10\}$, $L_0 = 360$, $N = 3960$, $k = 540$.
Range: $[360, 3600]$.

\medskip
\begin{tabular}{rrl}
\toprule
$e$ & $e \cdot k \bmod N$ & In range? \\
\midrule
1  & 540  & \checkmark \\
2  & 1080 & \checkmark \\
3  & 1620 & \checkmark \\
4  & 2160 & \checkmark \\
5  & 2700 & \checkmark \\
6  & 3240 & \checkmark \\
8  & 360  & \checkmark \\
9  & 900  & \checkmark \\
10 & 1440 & \checkmark \\
11 & 1980 & \checkmark \\
\bottomrule
\end{tabular}

\bigskip
\noindent\textbf{Omit 8.}
$F = \{2,3,4,5,6,7,9,10\}$, $L_0 = 1260$, $N = 13860$, $k = 1680$.
Range: $[1260, 12600]$.

\medskip
\begin{tabular}{rrl}
\toprule
$e$ & $e \cdot k \bmod N$ & In range? \\
\midrule
1  & 1680  & \checkmark \\
2  & 3360  & \checkmark \\
3  & 5040  & \checkmark \\
4  & 6720  & \checkmark \\
5  & 8400  & \checkmark \\
6  & 10080 & \checkmark \\
7  & 11760 & \checkmark \\
9  & 1260  & \checkmark \\
10 & 2940  & \checkmark \\
11 & 4620  & \checkmark \\
\bottomrule
\end{tabular}

\bigskip
\noindent\textbf{Omit 9.}
$F = \{2,3,4,5,6,7,8,10\}$, $L_0 = 840$, $N = 9240$, $k = 1008$.
Range: $[840, 8400]$.

\medskip
\begin{tabular}{rrl}
\toprule
$e$ & $e \cdot k \bmod N$ & In range? \\
\midrule
1  & 1008 & \checkmark \\
2  & 2016 & \checkmark \\
3  & 3024 & \checkmark \\
4  & 4032 & \checkmark \\
5  & 5040 & \checkmark \\
6  & 6048 & \checkmark \\
7  & 7056 & \checkmark \\
8  & 8064 & \checkmark \\
10 & 840  & \checkmark \\
11 & 1848 & \checkmark \\
\bottomrule
\end{tabular}

\bigskip
\noindent\textbf{Omit 10.}
$F = \{2,3,4,5,6,7,8,9\}$, $L_0 = 2520$, $N = 27720$, $k = 2750$.
Range: $[2520, 25200]$.

\medskip
\begin{tabular}{rrl}
\toprule
$e$ & $e \cdot k \bmod N$ & In range? \\
\midrule
1  & 2750  & \checkmark \\
2  & 5500  & \checkmark \\
3  & 8250  & \checkmark \\
4  & 11000 & \checkmark \\
5  & 13750 & \checkmark \\
6  & 16500 & \checkmark \\
7  & 19250 & \checkmark \\
8  & 22000 & \checkmark \\
9  & 24750 & \checkmark \\
11 & 2530  & \checkmark \\
\bottomrule
\end{tabular}

\medskip
\noindent All 9 templates verified. $\square$

%% ===================================================================
\begin{thebibliography}{12}

\bibitem{Wil67}
  J.~M.~Wills,
  \textit{Zwei S\"{a}tze \"{u}ber inhomogene diophantische Approximation
    von Irrationalzahlen},
  Monatsh.\ Math.\ \textbf{71} (1967), 263--269.

\bibitem{Cus73}
  T.~W.~Cusick,
  \textit{View-obstruction problems},
  Aequationes Math.\ \textbf{9} (1973), 165--170.

\bibitem{BGG+98}
  W.~Bienia, L.~Goddyn, P.~Gvozdjak, A.~Seb{\H{o}}, and M.~Tarsi,
  \textit{Flows, view obstructions, and the lonely runner},
  J.\ Combin.\ Theory Ser.\ B \textbf{72} (1998), no.~1, 1--9.

\bibitem{BH01}
  T.~Bohman and R.~Holzman,
  \textit{Six lonely runners},
  Electron.\ J.\ Combin.\ \textbf{8} (2001), no.~2, R3.

\bibitem{BS08}
  J.~Barajas and O.~Serra,
  \textit{The lonely runner with seven runners},
  Electron.\ J.\ Combin.\ \textbf{15} (2008), no.~1, R48.

\bibitem{Tao18}
  T.~Tao,
  \textit{Some remarks on the lonely runner conjecture},
  Contrib.\ Discrete Math.\ \textbf{13} (2018), no.~2, 1--31.

\bibitem{Ros25}
  N.~Rosenfeld,
  \textit{The Lonely Runner Conjecture for Eight Runners},
  arXiv:2509.14111 [math.CO], 2025.

\bibitem{Tra25}
  N.~Trakulthongchai,
  \textit{The Lonely Runner Conjecture for 9 and 10 Runners},
  arXiv:2511.22427 [math.CO], 2025.

\bibitem{Pal26}
  T.~Palelei,
  \textit{A Proof of the Lonely Runner Conjecture},
  Preprint, 2026.
  [Context reference; no proof in this paper depends on~\cite{Pal26}.]

\end{thebibliography}

\end{document}
